\section{The segment forest and its atoms}\label{dir:forest-definition}
The direct graph arguments below use the same finite-height segment forest,
with James's disjoint-segment norm \cite[p.~739]{James1974}.
We first give its norm and the two forms of its atomic unit ball that the
arguments need.

Let
\[
 T=\coprod_{n\ge1} T_n,\qquad T_n=\mathbb N^{\le n},
\]
where each component has its own root and is ordered by extension.  A segment
is a set $[a,b]=\{t:a\le t\le b\}$ with $a\le b$ in one component.
Disjoint segments need only be disjoint as sets. A set is ancestral if it contains every ancestor of each of its vertices.  On real finitely supported
arrays define
\begin{equation}\label{dir:atomic:forest-norm}
 \rho(y)=
 \sup_{\mathcal I}
 \left(\sum_{I\in\mathcal I}\left|\sum_{t\in I}y(t)\right|^2\right)^{1/2},
\end{equation}
where the supremum runs over finite disjoint families of segments.  Let $J$
be its completion and $X=J^*$.  Write $e_t^*$ for the coordinate functional
and $s_I^*=\sum_{t\in I}e_t^*$, with $s_\varnothing^*=0$.  We recall the elementary facts that the
argument needs, including the exact form of the atomic representation.

\begin{lemma}\label{dir:atomic:atoms}
The spaces $J$ and $X$ are reflexive, and finite coordinate arrays are dense
in both.  If a set $B\subset T$ meets each segment in either a segment or the
empty set, its coordinate projection $P_B$ is contractive on both spaces.
For an antichain $V$, the coordinate space supported on $V$ has its
$\ell_2(V)$ norm in both spaces.  In particular,
\begin{equation}\label{dir:atomic:atom-bound}
 \left\|\sum_{I\in\mathcal I}\lambda_I s_I^*\right\|_X
 \le\left(\sum_{I\in\mathcal I}\lambda_I^2\right)^{1/2}
\end{equation}
for every disjoint segment family.
Call a vector on the left with $\sum_I\lambda_I^2\le1$ an atom.  The unit
ball of $X$ is the norm-closed convex hull of the atoms.  If $S$ is finite
and ancestral, the unit ball of $P_SX$ is the convex hull of atoms
whose segments lie in $S$; every point has a finite convex representation.
\end{lemma}

\begin{proof}
Singleton tests give the lower estimate in
\[
 \|y\|_2\le\rho(y)\le\sqrt{n+1}\,\|y\|_2
 \qquad(\operatorname{supp}y\subset T_n).
\]
For the upper estimate apply Cauchy--Schwarz on every segment and use its
length at most $n+1$, followed by disjointness.  Norm squares add over the
components: one inequality follows by restricting each segment family to
each component, and the reverse inequality follows by combining families
that approximate the component suprema.  Thus $J$ is the $\ell_2$-sum of
spaces isomorphic to $\ell_2(T_n)$.  The dual of an $\ell_2$-sum is the
$\ell_2$-sum of the duals: Cauchy--Schwarz gives the upper norm bound and
finite coordinate tests give the reverse bound.  Applying this description
to the bidual proves reflexivity, and truncation first in components and
then in coordinates proves density in $J$ and $X$.

Intersecting a disjoint segment family with $B$ proves contractivity on
$J$, hence on $X$ by duality.  This applies to ancestral sets, their
complements, differences of nested ancestral sets, and antichains.  Each
segment meets an antichain in at most one point, so singleton tests identify
its $\rho$ norm with $\ell_2$; duality and the contractive projection give the
same identification in $X$.  Equation~\eqref{dir:atomic:atom-bound} follows
directly from Cauchy--Schwarz and \eqref{dir:atomic:forest-norm}.

The supremum of pairing with the atoms equals the right side of
\eqref{dir:atomic:forest-norm}.  If their norm-closed convex hull omitted a
point of the unit ball, Hahn--Banach separation would give a separating
functional in $X^*=J$, contradicting this equality.  Here reflexivity is
used to identify the separating functional with an element of $J$.  For
finite ancestral $S$, the same argument is finite-dimensional: projection
clips atoms to atoms in $S$, and the family of atoms in $S$ is compact,
being a finite union of continuous images of finite-dimensional Euclidean
unit balls.  Its convex hull is compact and equals the coordinate unit
ball.  Every point therefore has a finite convex representation.
\end{proof}

For ancestral $S$ put $Q_S=I-P_S$.  Finite-support density and
contractivity imply that, for a finite collection of vectors and any
$\eta>0$, a finite ancestral $S$ makes every $Q_S$-tail smaller than
$\eta$.  Indeed, take finite approximations and add all ancestors of their
supports; each vertex has only finitely many ancestors.


\section{Direct descent with primal tests}\label{dir:accumulated}

The next two arguments follow actual edges of a diamond and construct a
test vector at each step. Earlier diamond obstructions include
Lee--Naor's telescoping squared-distance inequality
\cite[Lemma~2.2]{LeeNaor2004} and Johnson--Schechtman's iterated
modulus-of-convexity estimate \cite[Section~4]{JohnsonSchechtman2009}.
Here the accumulated evaluation is compared with the cost of its primal
test. The first argument pays a bounded squared-norm cost for
each increase in its evaluation.  The second balances prefix sums of
comparable sizes and controls the final test all at once.

\subsection{The graph slopes}\label{dir:graph-slopes}
For an oriented parent edge of represented length $\ell$, use the
slope $x$ and deviations $z_j$ of Section~\ref{met:graph-definition}:
\begin{equation}\label{dir:splitting}
 x=\frac{f(b)-f(a)}\ell,\qquad
 z_j=\frac{2f(m_j)-f(a)-f(b)}\ell,\qquad
 \text{child slopes }x+z_j,\ x-z_j.
\end{equation}
In the noncontracting normalization all slopes have norm at most $D$
and the deviations are $2$-separated. After division by $D$, these
bounds are one and $2/D$. The symmetric splitting is the finite-level
version of the sprawling-tree construction in
\cite[Definition~3.7 and Proposition~6.4]{BassetLancienProchazka2026}.

The selections below use the same compactness step as
Lemma~\ref{met:finite-pair}: two bounded projections onto a fixed finite
head can be made arbitrarily close. Their separated full vectors then
have a large tail difference, so at least one has a large tail. This
choice requires the infinitely many middle vertices. Once a middle
vertex has been selected, either sign in~\eqref{dir:splitting} follows
an actual child edge.

\subsection{A correction at the first fresh roots}

The correction controls the extra squared cost of segments crossing
the head boundary. Its value is placed at each first exterior vertex.
For crossings with a large head sum, the proof either makes the cross
term nonpositive or controls it by splitting the head portion.

\begin{lemma}\label{dir:root-correction}
Suppose $v,u\in c_{00}(T)$, $v=P_Av$, and $u=Q_Au$, where $A$ is finite
and ancestral.  There is $c=Q_Ac\in c_{00}(T)$ such that
\begin{equation}\label{dir:correction-cost}
 \rho(c)\le\rho(u),\qquad
 \rho(v+c+r u)^2\le\rho(v)^2+16\rho(u)^2
 \quad(r\in\{-1,1\}).
\end{equation}
\end{lemma}
\begin{proof}
The complement of $A$ is the disjoint union of subtrees $U_d$ whose
roots are its first vertices.  Put $u_d=u|_{U_d}$ and $a_d=\rho(u_d)$.
Independent segment tests in these subtrees give
$\sum_d a_d^2=\rho(u)^2$; only finitely many terms are nonzero.
If $d$ has parent $p\in A$, examine the suffix sums
$V_d(h)=\sum_{q\in[h,p]}v(q)$.  If some suffix has absolute sum greater
than $2a_d$, choose the shortest such suffix, with initial vertex $h_d$, and set
\[
 c_d=-\operatorname{sign}(V_d(h_d))a_d.
\]
Otherwise set $c_d=0$, also when $d$ has no parent in $A$.  Let
$c=\sum_dc_de_d$.  The roots $d$ are incomparable, so
$\rho(c)^2=\sum_dc_d^2\le\rho(u)^2$.

Fix a disjoint segment family testing $v+c+ru$.  Segments wholly outside
$A$ cost at most $4\sum_da_d^2$ in squared sum, since
$\rho(c_de_d+ru_d)\le2a_d$.  A crossing into $U_d$ has head sum $P$ and
tail sum $Q$, with $|Q|\le2a_d$.  We claim that
\[
 (P+Q)^2\le E+12a_d^2,
\]
where $E$ is the squared-sum test of one or two disjoint subsegments of
the original head portion.  If $|P|\le2a_d$, use $E=P^2$, since
$2|PQ|+Q^2\le12a_d^2$.  Otherwise the preferred suffix exists.  If its
sum has the same sign as $P$, the choice of $c_d$ gives $PQ\le0$, and
again $E=P^2$ suffices.  In the remaining case let $P'=V_d(h_d)$.
The preferred suffix lies properly inside the head portion: the latter
qualifies, and the former is shortest and has the opposite sign.
Splitting there gives
\[
 E=(P-P')^2+(P')^2\ge P^2+4|P|a_d\ge P^2+2PQ.
\]
Adding $Q^2\le4a_d^2$ proves the claim in this case as well.
All the replacement head segments remain disjoint.  Their total squared
test, including the unchanged segments in $A$, is at most $\rho(v)^2$.
At most one crossing uses each fresh root $d$.  Combining the crossing
cost $12\sum_da_d^2$ and the wholly exterior cost proves
\eqref{dir:correction-cost}.
\end{proof}

\begin{proposition}\label{dir:deterministic-descent}
Every embedding of $\mathcal D_k$ into $X$ with distortion at most $D$
satisfies $k\le256D^2$.
\end{proposition}
\begin{proof}
Rescale to upper Lipschitz constant one and lower constant $1/D$.
Start at the initial edge with test $v=0$.  At a current edge of slope
$x$, enlarge a finite ancestral $A$ to contain the support of $v$ and
make $\|Q_Ax\|<1/(4D)$.  In \eqref{dir:splitting}, choose two deviations
whose head projections are less than $1/(2D)$ apart.  Since the
deviations are $2/D$-separated, at least one of them, denoted $z$, has
$\|Q_Az\|>1/(2D)$.  Finite-support density in $J$ and contractivity of
$Q_A$ supply $u=Q_Au\in c_{00}(T)$ with
$\rho(u)\le1$ and $z(u)>1/(2D)$.

Apply Lemma~\ref{dir:root-correction}.  Because $\rho(c)\le1$,
$|x(c)|<1/(4D)$.  Averaging the two matching evaluations gives
\[
 \frac12\sum_{r=\pm1}(x+rz)(v+c+ru)
 =x(v)+x(c)+z(u)>x(v)+\frac1{4D}.
\]
Choose a sign attaining at least this average, follow that actual child
edge, and replace $v$ by $v+c+ru$.  Each step increases $\rho(v)^2$ by
at most $16$.  After $k$ steps the evaluation is at least $k/(4D)$,
whereas the final slope has norm at most one and the final test has
norm at most $4\sqrt{k}$.  Hence $k\le256D^2$.
\end{proof}

\subsection{Balancing dyadic sizes along a random descent}
The next construction does not control the increase of the test norm
at an individual step.  It controls prefix sums separately at every
dyadic size, then bounds any segment by charging its largest size to a
root through which that segment passes.

\begin{proposition}\label{dir:bin-descent}
Every noncontracting $D$-Lipschitz embedding of $\mathcal D_k$ into $X$
satisfies $D>\sqrt{k}/144$ for $k\ge1$.
\end{proposition}
\begin{proof}
We descend at random, choosing a fair sign after selecting a middle
vertex.  Write $x$ for the current slope and $G$ for the accumulated
test, initially zero.  Along a fixed history, the preceding tests lie
in finite blocks $Q\setminus P$ with ancestral endpoints, and every
new $P$ contains the previous $Q$.  Enlarge $P$ until
$\|Q_Px\|<1/32$.  By the separation in \eqref{dir:splitting} and
finite-dimensional total boundedness, choose a deviation $z$ with
$\|Q_Pz\|\ge3/4$.  There is a finite test $g$ supported in a block
$Q\setminus P$ such that
\begin{equation}\label{dir:bin-witness}
 \rho(g)\le1,\qquad z(g)>1/2.
\end{equation}
Here $Q$ is a finite ancestral enlargement containing the support of
$g$; both choices are made before tossing the next sign.

Decompose the block into its connected tree components $C$.  If $r_C$
is the root of $C$, define
\[
 p_C(v)=\sum_{w\in[r_C,v]}g(w),\qquad
 A_C=\max_{v\in C}|p_C(v)|.
\]
Ignore components with $A_C=0$.  Choose a power of two $B_C$ with
$A_C\le B_C<2A_C$.  Disjoint root-prefix tests give
$\sum_CA_C^2\le1$.  Differencing the potentials
$\max(p_C,0)$ and $\min(p_C,0)$ gives arrays $g_C^+,g_C^-$ with
$g|_C=g_C^++g_C^-$.  On every subsegment, either array has sum in
absolute value at most the absolute value of the corresponding sum of
$g$, since both scalar truncations are 1-Lipschitz.

To state the balancing rule precisely, call a component earlier if it
was selected before the current update. Label each earlier component by
its bin size and, for every vertex $v$ and dyadic size $B$, put
\[
 S_B(v)=\sum_{\substack{t\preceq v:\ t\text{ lies in an earlier}\\
                              \text{component of bin size }B}}G(t).
\]
The invariant is $|S_B(v)|\le B$ for every $v$. Before adding a new
component $C$, the sum of earlier bin-$B_C$ contributions above its root
$r_C$ is fixed; choose a weak sign opposite to this sum, with either
choice at zero. Allow the new test on $C$ to be zero or either of
$g_C^+,g_C^-$ oriented so that all its prefix sums have this chosen sign.
Each such prefix has magnitude at most $B_C$.

This preserves the invariant at every vertex. All earlier supports lie
in $P$, so their prefix contribution is constant through $C$. Adding
an opposite-sign quantity of magnitude at most $B_C$ keeps its absolute
value at most $B_C$. In a gap there is no new contribution; an ignored
component has $g|_C=0$. Beyond $Q$, the new prefix is the one at its last
ancestor in $Q$, or zero when there is no such ancestor. Thus the same
bound holds after every update, for every ancestor path.

The sign restriction is fixed before the next fair sign $s\in\{-1,1\}$.
For each $C$, choose an allowed option maximizing $s z$ on that
component; let $w$ be the resulting increment.  The choice depends on
$s$.  Nevertheless $\rho(w)\le1$: every tested segment meets this block
in a segment in one component, and the preceding contraction estimate
applies simultaneously to disjoint tests.  If $a,b$ are the evaluations
under $z$ of the two oriented offers in a component, then
\[
 \mathbb E_s\max(0,sa,sb)
 \ge\frac{|a|+|b|}{4}
 =\frac{|z(g_C^+)|+|z(g_C^-)|}{4}
 \ge\frac{z(g|_C)}4.
\]
Sum over $C$ and use \eqref{dir:bin-witness} to obtain
$\mathbb E_s[s z(w)]>1/8$.  Also $|x(w)|<1/32$ and
$\mathbb E_s[s z(G)]=0$.  Thus the next actual slope $x+sz$ satisfies
\[
 \mathbb E_s[(x+sz)(G+w)]>x(G)+1/8-1/32>x(G)+1/16.
\]
After $k$ steps the expected evaluation exceeds $k/16$.  The entire
probability space has only $2^k$ outcomes.

It remains to bound $\rho(G)$ for every final outcome.  Fix any finite
disjoint segment family.  On a segment, first separate the portion
lying in the component that contains its initial vertex, if this is
a nonignored component.  Across all segments these initial portions
have squared-sum total at most $k$: grouping them by construction step
gives disjoint tests of an increment of norm at most one.
The remainder of a segment enters every nonignored component it meets
at its root, since all blocks have ancestral endpoints and are
successively outside previous blocks.  For a bin of size $B$, its sum
on this remainder is a difference of two values of $S_B$, hence
has magnitude at most $2B$. A remainder meeting no nonignored component
contributes zero. Otherwise, if $B_*$ is the largest encountered bin,
geometric summation gives a total magnitude at most $4B_*$.  Charge the
square of this bound to an encountered component with size $B_*$.
Two disjoint segments cannot charge the same component: both would
contain its root.  Since
\[
 \sum_{\text{all components}}B_C^2
 \le4\sum_{\text{steps}}\sum_CA_C^2\le4k,
\]
the remaining portions have squared-sum total at most $64k$.
The Euclidean triangle inequality for the vector of segment sums now
gives $\rho(G)\le\sqrt{k}+8\sqrt{k}=9\sqrt{k}$.

Every final slope has norm at most $D$, so every final evaluation is at
most $9D\sqrt{k}$.  Comparison with its expectation gives
$k/16<9D\sqrt{k}$, as claimed.
\end{proof}

\section{Direct graph obstructions from atomic prediction energy}
\label{dir:atomic:section}

We give two direct proofs that the countably branching diamonds do not embed
with uniformly bounded distortion into the finite-height segment-forest dual.
Both proofs first select finitely many edge slopes from the infinitely
branching graph.  An atomic representation then turns these slopes into
conditional means.  The first proof counts segment endpoints and predicts the
coefficients of segments crossing a long interval of layers.  The second
counts segment starts and follows each segment until it terminates.  Their
energy estimates concern these finite martingales; no midpoint-modulus
estimate is needed.

\begin{theorem}\label{dir:atomic:obstruction}
For the real space $X$ defined by \eqref{dir:atomic:forest-norm}, no finite
$D$ has the following property: for every $k\ge1$ there are a map
$f_k:\mathcal D_k\to X$ and a scalar $a_k>0$ such that
\[
 a_k d_k(u,v)\le\|f_k(u)-f_k(v)\|_X\le D a_k d_k(u,v)
 \qquad(u,v\in\mathcal D_k).
\]
\end{theorem}

Rescaling by $D a_k$ reduces the theorem to excluding maps satisfying
\begin{equation}\label{dir:atomic:normalization}
 D^{-1}d_k(u,v)\le\|f(u)-f(v)\|_X\le d_k(u,v).
\end{equation}
Necessarily $D\ge1$. The graph normalization and old-vertex distance
calculation in Section~\ref{met:graph-definition} apply throughout.
Orient the edges consistently and write their normalized slopes as in
\eqref{met:slopes}. They have norm at most one, and each retained parent
has the symmetric children
\begin{equation}\label{dir:atomic:splittings}
 x+d_j,\quad x-d_j,\qquad
 \|d_j-d_l\|_X\ge 2/D\quad(j\ne l).
\end{equation}

At every retained edge choose one replacement path, retain both its
slopes, and then descend by fair binary choices.  The resulting random
slopes $M_0,\ldots,M_k$ form a martingale, because the two children average
to their parent.  Let $\mathcal F_i$ record the first $i$ binary choices.
Every history at level $i$ has probability $2^{-i}$, and each history has
two children of conditional probability $1/2$.  Infinite branching is
used only to select a suitable replacement path before this finite
probability space is formed.

\begin{lemma}[A finite slope martingale and its exact atomic law]
\label{dir:atomic:preparation}
Suppose a depth-$k$ embedding satisfies~\eqref{dir:atomic:normalization}.
For every $0<\eta<1/(4D)$ one can select a fair binary slope martingale
$(M_i)_{i=0}^k$, nested finite ancestral sets $S_0\subseteq\cdots\subseteq S_k$,
and a finitely valued random atom $Z$ supported in $S_k$, such that, with
$E_i=\mathbb E[\,\cdot\mid\mathcal F_i]$ and $U_i=E_iZ=P_{S_k}M_i$,
\begin{equation}\label{dir:atomic:preparation-bounds}
 \|Q_{S_{i-1}}U_{i-1}\|<\eta,\qquad
 \|P_{S_i\setminus S_{i-1}}U_i\|\ge1/D-4\eta
 \quad(1\le i\le k).
\end{equation}
The unprojected slopes also satisfy $\|Q_{S_i}M_i\|<\eta$ for
$0\le i\le k$. Every $M_i$ and $U_i$ has norm at most one. The additional terminal
sampling used to produce $Z$ is excluded from each $\mathcal F_i$,
including $\mathcal F_k$.
\end{lemma}
\begin{proof}
Choose a finite ancestral $S_0$ with $\|Q_{S_0}M_0\|<\eta$.
Suppose that the selected level-$(i-1)$ slopes have tails smaller than
$\eta$ outside a common finite ancestral $S_{i-1}$.  For each such parent
$x$, the first-child slopes $x+d_j$ in \eqref{dir:atomic:splittings} have
bounded projections to the finite-dimensional space $P_{S_{i-1}}X$.
Two distinct projections are within $2\eta$.  Their tail difference has
norm at least $2/D-2\eta$, so one first-child tail has norm at least
$1/D-\eta$.  Retain that path.  Since its two slopes add to $2x$, the
mate's tail has norm at least $1/D-3\eta$.

Only finitely many slopes have now been selected.  Enlarge $S_{i-1}$ to a
finite ancestral $S_i$ making all their tails smaller than $\eta$.
Writing $B_i=S_i\setminus S_{i-1}$, both children satisfy
\[
 \|P_{B_i}M_i\|\ge1/D-4\eta.
\]
Repeat through level $k$.  All sets are common to an entire level, not
chosen after its final random history.

The finite-support martingale $P_{S_k}M_i$ has norm at most one.
By Lemma~\ref{dir:atomic:atoms}, at every terminal binary history its value
is a finite convex combination of atoms supported in $S_k$.  Enlarge the
finite sample space by selecting one atom from that combination and call
it
\[
 Z=\sum_{I\in\mathcal I}a_I s_I^*,\qquad \sum_{I\in\mathcal I}a_I^2\le1.
\]
The extra terminal sampling is excluded from every $\mathcal F_i$, even
$\mathcal F_k$.  With $E_i=\mathbb E[\,\cdot\mid\mathcal F_i]$, the tower
property gives $E_iZ=P_{S_k}M_i$. The projections commute, and
$S_i\subseteq S_k$, so the displayed
layer lower bounds also hold for $U_i$. Moreover,
$Q_{S_{i-1}}U_{i-1}=P_{S_k}Q_{S_{i-1}}M_{i-1}$; contractivity gives
the first bound. This proves~\eqref{dir:atomic:preparation-bounds}.
\end{proof}

The lemma supplies the same finite martingale for two different energy
arguments. In either argument a large child layer is decomposed into
three terms: segments charged to the chosen interval, coefficients
already predicted from the parent, and the change in their prediction.
The first term has a counting budget, the second is controlled by the
small parent tail, and scalar $L^2$ orthogonality bounds the third.

\subsection{Endpoint intervals and crossing predictions}
\label{dir:atomic:endpoint-route}

We prove Theorem~\ref{dir:atomic:obstruction} first by charging the two
endpoints of every terminal atom.  Assume that \eqref{dir:atomic:normalization}
is available for every $k$, with $D$ fixed, and set
\begin{equation}\label{dir:atomic:endpoint-constants}
 \eta=\frac1{32D},\qquad c=\frac1{2D}.
\end{equation}
Apply Lemma~\ref{dir:atomic:preparation} with this $\eta$, write
$B_i=S_i\setminus S_{i-1}$, and write its random atom as
\[
 Z=\sum_{I\in\mathcal I}w_I s_I^*,\qquad \sum_Iw_I^2\le1.
\]
Since $1/D-4\eta>c$, we have on every binary history
\begin{equation}\label{dir:atomic:endpoint-layer}
 \|Q_{S_{i-1}}E_{i-1}Z\|<\eta,\qquad
 \|P_{B_i}E_iZ\|\ge c
 \quad(1\le i\le k).
\end{equation}

Take $k=ML$ and partition the layers into $M$ consecutive intervals of
length $L$.  A segment has at most two endpoints, so the sum, over these
intervals, of its squared coefficient charged to an interval containing
an endpoint is at most $2w_I^2$.  Consequently some interval $p\le i\le q$,
$q-p+1=L$, has
\[
 \mathbb E\sum_{\substack{I:\text{an endpoint of }I\text{ lies in}
                   \,S_q\setminus S_{p-1}}}w_I^2\le\frac2M.
\]
Call these segments endpoint segments and let $Z^{\rm end}$ be their
sum.  Put $b_0=\sqrt{2/M}$.  By \eqref{dir:atomic:atom-bound},
\begin{equation}\label{dir:atomic:endpoint-error}
 \|Z^{\rm end}\|_{L^2(X)}\le b_0.
\end{equation}
Here $\|V\|_{L^2(X)}=(\mathbb E\|V\|_X^2)^{1/2}$; contractive
projections and conditional expectations preserve upper bounds in this
norm, by Jensen's inequality.

Every other segment meeting a layer in this interval begins in
$S_{p-1}$ and ends outside $S_q$: heredity puts its first endpoint in
$S_q$, and the absence of either endpoint in $S_q\setminus S_{p-1}$
then forces these two alternatives.  Call such segments crossing
segments and write $Z^{\rm cr}$ for their sum.  On every relevant layer
only $Z^{\rm end}$ and $Z^{\rm cr}$ contribute.

For $p\le i\le q+1$ let $A_i$ be the minimal vertices of
$S_k\setminus S_{i-1}$, and put $A_i'=A_i\cap S_i$ for $i\le q$.
These are finite antichains.  Each crossing segment meets every $A_i$
in exactly one vertex, since it begins in $S_{p-1}$ and ends outside
$S_q$.  Label it by its vertex $r\in A_p$.  Disjointness allows at most
one crossing segment with a given label.  Let $W_r$ be its coefficient,
and $d_i^r$ its vertex in $A_i$; in the absence of such a segment put
$W_r=0$ and $d_i^r=\dagger$ for every $i$.

A later location determines an earlier one by taking its unique
$A_i$-ancestor, with $\dagger$ mapping to $\dagger$.  For each fixed
$r$, the sigma-algebras
\[
 \mathcal G_{i-1}^r=\mathcal F_{i-1}\vee\sigma(d_i^r)
 \qquad(p\le i\le q+1)
\]
therefore increase with $i$.  Define $H_{i-1}^r=\mathbb E[W_r\mid
\mathcal G_{i-1}^r]$.  Scalar conditional expectation is orthogonal
projection in $L^2$, whence
\begin{equation}\label{dir:atomic:crossing-energy}
 \sum_{i=p}^q h_i^2\le1,\qquad
 h_i^2:=\sum_{r\in A_p}\mathbb E|H_i^r-H_{i-1}^r|^2.
\end{equation}
Indeed the sum for each $r$ telescopes to at most $\mathbb E W_r^2$,
and the sum of those quantities is at most one.  Different labels may
use different sigma-algebras; no joint refinement across labels is
being asserted.

Let $J_i^r$ be the part of the crossing segment with label $r$ in $B_i$,
empty if the label is absent.  It is determined by $d_{i+1}^r$: the
previous location is recovered from it, and the portion runs from
$d_i^r$ up to but excluding $d_{i+1}^r$.  If these locations agree the
portion is empty.  Therefore
\[
 E_i\sum_r W_r s_{J_i^r}^*=E_i\sum_r H_i^r s_{J_i^r}^*.
\]
Replacing $H_i^r$ by $H_{i-1}^r$ in this expression costs at most $h_i$
in $L^2(X)$, by disjointness, \eqref{dir:atomic:atom-bound}, and
conditional Jensen.  Together with \eqref{dir:atomic:endpoint-error},
this reduces the layer estimate to bounding
\[
 V_i=E_i\sum_r H_{i-1}^r s_{J_i^r}^*.
\]

Fix a history through time $i-1$.  For $v\in A_i'$, let $r(v)$ be its
unique $A_p$-ancestor and let $m(v)$ be the mean of $W_{r(v)}$
conditional on this history and $d_i^{r(v)}=v$.  Set it to zero on an
empty conditional event.  Write
\[
 p_j(v)=\Pr(d_i^{r(v)}=v\mid\mathcal F_j),\qquad j=i-1,i.
\]
On either child history, $p_i(v)\le2p_{i-1}(v)$, because that history
has half its parent's probability.  The contribution from $v$ to $V_i$
is $m(v)p_i(v)$ times an average of segment functionals starting at $v$.
The descendant subtrees at different $v$ are disjoint.  Independently
sampling one segment from each of these averages and applying
\eqref{dir:atomic:atom-bound} before averaging proves
\[
 \|V_i\|_X\le
 \left(\sum_{v\in A_i'}|m(v)p_i(v)|^2\right)^{1/2}
 \le2\left(\sum_{v\in A_i'}|m(v)p_{i-1}(v)|^2\right)^{1/2}.
\]
The products in the final sum are precisely the coordinates of
$E_{i-1}Z^{\rm cr}$ on $A_i'$.  The antichain norm, the parent-tail
bound in \eqref{dir:atomic:endpoint-layer}, and
\eqref{dir:atomic:endpoint-error} now give
\[
 \|V_i\|_{L^2(X)}\le2(\eta+b_0).
\]
Combining the three contributions yields, for every $p\le i\le q$,
\begin{equation}\label{dir:atomic:endpoint-final}
 c\le2\eta+3\sqrt{2/M}+h_i.
\end{equation}
By \eqref{dir:atomic:crossing-energy}, some $h_i\le L^{-1/2}$.
Since $2\eta=c/8<c$, choose $M$ and then $L$ so that
$3\sqrt{2/M}+L^{-1/2}<c-2\eta$.  This contradicts
\eqref{dir:atomic:endpoint-final} and proves
Theorem~\ref{dir:atomic:obstruction}.

\subsection{Start intervals and terminating segments}
\label{dir:atomic:status-route}

A second proof uses a smaller counting budget, one charge per segment,
and a different prediction space.  It allows old segments to terminate
inside the chosen interval.  Recording their last vertices is the
additional information needed to retain a refining sequence of
partitions.

Assume again that \eqref{dir:atomic:normalization} is available for all
$k$ at a fixed $D$.  This time choose
\begin{equation}\label{dir:atomic:status-constants}
 c_0=\frac1{2D},\qquad \delta=\frac{c_0}{2},\qquad
 \eta>0,\quad3\eta<\frac{c_0}{2},\quad4\eta<\frac\delta4.
\end{equation}
Apply Lemma~\ref{dir:atomic:preparation}, put $H_i=S_i$ and
$J_i=H_{i+1}\setminus H_i$, and write
\[
 Z=\sum_{I\in\mathcal I}\lambda_I s_I^*,\qquad
 \sum_I\lambda_I^2\le1,\qquad U_i=E_iZ=P_{S_k}M_i.
\]
The common preparation gives the exact conditional-mean identity.
Since $1/D=2c_0$ and $\eta<c_0$, its bounds imply
\begin{equation}\label{dir:atomic:status-selected}
 \|Q_{H_i}M_i\|<\eta,\qquad
 \|P_{J_i}M_{i+1}\|\ge2c_0-4\eta,
\end{equation}
and, in particular, the weaker bounds
\begin{equation}\label{dir:atomic:status-layer}
 \|Q_{H_i}U_i\|\le2\eta,\qquad
 \|P_{J_i}U_{i+1}\|\ge c_0-3\eta>\delta.
\end{equation}
We keep this slack in the estimates below. The distinction from the
endpoint argument will be that segments may terminate inside the chosen
interval; their last vertices must therefore remain in the prediction
information.

Take $k=NL$ and divide its steps into intervals $r\le i<r+L$, where
$r$ ranges over $0,L,\ldots,(N-1)L$.  For each interval put
\[
 B_r=\sum_{\min I\in H_{r+L}\setminus H_r}\lambda_I^2,
 \qquad b_r=\mathbb E B_r.
\]
Every segment has one first vertex, so $\sum_r b_r\le1$.  Fix one such
interval.  Call a sampled segment old if its first vertex lies in $H_r$,
and new if its first vertex lies in $H_{r+L}\setminus H_r$.
No other segment meets a layer $J_i$ of this interval, by heredity.
The new contribution after projection and conditional expectation has
$L^2(X)$ norm at most $\sqrt{b_r}$, because its pointwise norm before
conditioning is at most $\sqrt{B_r}$.

To predict the coefficients of old segments, use the finite set of pairs
$(\omega,I)$ with $\omega$ a terminal sample outcome and $I$ an old
segment in its sampled atom.  Give each such pair mass $\Pr(\omega)$,
obtaining a finite measure $\mu$, and write
$\lambda(\omega,I)=\lambda_I(\omega)$.  Then
\[
 \int\lambda^2\,d\mu
 =\mathbb E\sum_{I\text{ old}}\lambda_I^2\le1.
\]
If there are no pairs, all old contributions and all energies below
are zero.  Otherwise, for $r\le i\le r+L$, partition this pair set by
the binary history through time $i$ and the following tagged status:
\begin{itemize}
 \item if $I$ has a vertex outside $H_i$, record ``active'' and its
       first vertex outside $H_i$;
 \item if $I\subset H_i$, record ``inactive'' and its last vertex.
\end{itemize}
These partitions refine.  An old segment starts in $H_r\subset H_i$,
so a later active entry determines its earlier entry by the ancestor
chain.  A later inactive endpoint determines whether the earlier status
was inactive at that endpoint or active at the first ancestor outside
$H_i$.  The tags distinguish these cases.

Let $m_i$ be the mean of $\lambda$ on each positive-mass cell of the
$i$th partition, with arbitrary value zero on empty cells, and define
\begin{equation}\label{dir:atomic:status-energy}
 g_i=\int|m_{i+1}-m_i|^2\,d\mu.
 \qquad
 \sum_{i=r}^{r+L-1}g_i\le\int\lambda^2\,d\mu\le1.
\end{equation}
The inequality is the orthogonality identity for conditional means in
$L^2(\mu)$; it does not require $\mu$ to be a probability measure.

The next status determines $I\cap J_i$.  If it is active, this portion
ends just before the next entry; if it is inactive, it ends at the
recorded last vertex.  Its first vertex, when the portion is nonempty,
is the first ancestor outside $H_i$.  Equal successive entries, or an
endpoint already in $H_i$, give an empty portion.  Thus the clipped
segment functional is constant on every cell of the $(i+1)$st
partition, and
\[
 E_{i+1}\sum_{I\text{ old}}\lambda_I P_{J_i}s_I^*
 =E_{i+1}\sum_{I\text{ old}}m_{i+1}(\omega,I)P_{J_i}s_I^*.
\]
Split $m_{i+1}=(m_{i+1}-m_i)+m_i$.  Disjointness of the clipped
segments and conditional Jensen bound the first term in $L^2(X)$ by
$\sqrt{g_i}$.

We now bound the second term using the small parent tail.  At a fixed
history $h$ through time $i$, possible active entries form the antichain
of minimal vertices outside $H_i$.  Only finitely many occur in the
finite sample space.  Let $V$ be those lying in $H_{i+1}$, let $p_v$
be the conditional expected number of old segments with entry $v$, and
let $m_v$ be their mean coefficient under $\mu$ conditional on $h$ and
this entry.  Empty cells give $m_v=0$.  There is at most one such segment
per outcome, by disjointness.  Hence the coordinates of the conditional
old contribution on $V$ are exactly $p_vm_v$.  Its only other
contributions in $U_i$ come from new segments.  The antichain norm and
\eqref{dir:atomic:status-layer} therefore imply
\begin{equation}\label{dir:atomic:status-antichain}
 \left(\sum_{v\in V}|p_vm_v|^2\right)^{1/2}
 \le2\eta+(E_iB_r)^{1/2}.
\end{equation}
For the new term, this follows from its disjoint-segment bound and
conditional Cauchy--Schwarz before restricting to $V$.

At a child history $h'$ of $h$, the conditional expected number of old
segments with a given previous entry $v$ is at most $2p_v$.  To see the
measure normalization explicitly, if $C(h,v)$ denotes the corresponding
cell of pairs, then
\[
 p_v=\frac{\mu(C(h,v))}{\Pr(h)},\qquad
 \frac{\mu(C(h,v)\cap\{\text{history }h'\})}{\Pr(h')}
 \le2p_v,
\]
since $\Pr(h')=\Pr(h)/2$.  This fairness assertion concerns the original
binary probability, not the pair measure $\mu$.

The contribution with coefficient $m_i$ at $h'$ is consequently, for
each $v\in V$, $m_v$ times a nonnegative mixture of segment functionals
starting at $v$, of total mass at most $2p_v$.  Add a zero functional to
make this mass exactly $2p_v$, when needed.  Choose independently one
term per $v$ from the normalized mixtures.  Distinct $v$ lie in disjoint
descendant subtrees, so \eqref{dir:atomic:atom-bound} bounds every such
sum by $2(\sum_v|p_vm_v|^2)^{1/2}$.  Averaging and then applying
\eqref{dir:atomic:status-antichain} gives an $L^2(X)$ bound
$4\eta+2\sqrt{b_r}$ for this term.  Adding the prediction difference
and the new contribution proves
\begin{equation}\label{dir:atomic:status-final}
 \delta\le4\eta+3\sqrt{b_r}+\sqrt{g_i}
 \qquad(r\le i<r+L).
\end{equation}

Choose $L>4/\delta^2$ first.  Then choose $N$ large enough that
$3/\sqrt N<\delta/4$, and take $k=NL$.  Some interval satisfies
$b_r\le1/N$, since $\sum_r b_r\le1$.  For this interval,
\eqref{dir:atomic:status-constants} and \eqref{dir:atomic:status-final}
force $\sqrt{g_i}>\delta/2$ at every step.  This contradicts
\eqref{dir:atomic:status-energy}, because
\[
 \sum_{i=r}^{r+L-1}g_i>L\delta^2/4>1.
\]
This completes the second proof of
Theorem~\ref{dir:atomic:obstruction}.

The two estimates retain different information.  In the endpoint proof,
the selected interval has endpoint energy at most $2/M$, all surviving
segments cross the entire interval, and the coefficient predictions are
scalar martingales for each boundary label.  In the start proof, total
start energy is at most one, segments may terminate inside the interval,
and prediction uses their refining active or inactive statuses on a
finite measure space.  These are direct obstruction arguments with
their respective parameter choices; neither calculation above asserts
an asymptotic midpoint modulus.

\section{Fresh coordinate layers along diamond slopes}\label{dir:layers}

We now use the linear path-cost space.  Its intrinsic estimates control
either a fresh part of a finite martingale or the squared-norm increase
of a symmetric split.  The graph arguments below supply precisely those
inputs, with their separate support and measurability conditions.

For $T_n=\mathbb N^{\le n}$, including its root, let
$b_t=\sum_{s\preceq t}e_s$ and
\[
 L_n(z)=\inf_{z=h+\sum_t\alpha_tb_t}
          \bigl(\|h\|_{\ell_2(T_n)}+\|\alpha\|_{\ell_1(T_n)}\bigr),
 \qquad
 X_{\rm lin}=\left(\bigoplus_{n\ge1}(\ell_2(T_n),L_n)\right)_2.
\]
The infimum ranges over real $h\in\ell_2(T_n)$ and
$\alpha\in\ell_1(T_n)$; it need not be attained.  The intrinsic estimates
below allow arbitrarily close decompositions.  This definition is the
linear path-cost model of
\cite[Propositions~5.1 and~5.2]{TreePotentials}.  In particular, that
paper proves the equivalence with the bounded global-test norm on each
finite-height component.  No infinite-height completion is substituted
for this Hilbert sum.

For a finite ancestral set $A$ in the forest, $P_A$ is contractive on
$X_{\rm lin}$: it restricts a Hilbert vector and truncates every $b_t$ to another
root path or zero.  Finite coordinate vectors are dense because
$\|z\|_2/\sqrt{n+1}\le L_n(z)\le\|z\|_2$ on each component.
Thus the $P_A$ approximate the identity strongly.  For $A\subseteq B$,
the layer projection is $P_{B\setminus A}=P_B-P_A$ and has norm at most
two.  The tail projection $I-P_A$ also has norm at most two.

We use three forms of coordinate selection. Exact graph slopes can be
paired with windows chosen at the current history, or with one head
fixed across a whole level. Estimates requiring exact support instead
use finite approximations: the approximating slopes retain the exact
parent--child averaging identity, and the next support is fixed before
the next sign. We give each intrinsic estimate immediately before the
construction that meets its hypotheses.

Throughout, the probability space is a uniform finite sign cube,
$\mathcal F_i$ records its first $i$ signs, and a dyadic martingale
satisfies $W_i=\mathbb E[W_m\mid\mathcal F_i]$. The estimates are the
separate sign-prediction, variation, and stopping arguments proved in
\cite{TreePotentials}; no constant is transferred between these proofs.

\subsection{Exact slopes with adaptive windows}\label{dir:adaptive-windows}
The weighted sign-prediction estimate
\cite[Theorem~6.4]{TreePotentials} supplies an absolute constant
$C_{\rm sign}$ such that, for finite ancestral windows
$A_i\subseteq B_i$ measurable with respect to $\mathcal F_i$, with
$B_i\subseteq A_{i+1}$ whenever $i+1<m$,
\begin{equation}\label{dir:sign-input}
 \sum_{i<m}\mathbb E\|P_{B_i\setminus A_i}W_{i+1}\|^2
 \le C_{\rm sign}\left(\mathbb E\|W_m\|^2+
       \sum_{i<m}\mathbb E\|P_{B_i\setminus A_i}W_i\|^2\right).
\end{equation}
There is no support assumption on $W_i$ and no specified numerical
value for $C_{\rm sign}$.

Suppose $f:\mathcal D_k\to X_{\rm lin}$ is noncontracting and $D$-Lipschitz, where
$k\ge1$.  We recursively select a middle vertex at each encountered
edge, then choose a half-edge by a fair sign.  The resulting true slopes
$W_i$ form a dyadic martingale and satisfy $\|W_i\|\le D$.
At a history of length $i$, choose a finite ancestral $A_i$, containing
$B_{i-1}$ when $i\ge1$, with
$\|(I-P_{A_i})W_i\|\le2^{-k}$.  For each middle vertex consider its two
child slopes.  The first-child sequence is pairwise $2$-separated, as
is the second-child sequence.  Each sequence has only finitely many
members with tail norm at most $1/2$: infinitely many such members,
together with compactness of their bounded head projections, would
give a pair at distance less than $2$.  Thus one middle vertex has
\emph{both} child tails greater than $1/2$.

Before tossing the next sign, enlarge $A_i$ to a finite ancestral $B_i$
capturing these two child tails closely enough that
\[
 \|P_{B_i\setminus A_i}W_{i+1}\|>1/4
\]
on either outcome.  For the parent,
$P_{B_i\setminus A_i}W_i=P_{B_i}(I-P_{A_i})W_i$, so its layer norm is
at most $2^{-k}$.  All windows are known before the next sign and are
disjoint along each complete history.  Applying
\eqref{dir:sign-input} therefore gives
\begin{equation}\label{dir:adaptive-conclusion}
 \frac{k}{16}\le C_{\rm sign}(D^2+k2^{-2k})
 \le C_{\rm sign}(D^2+1).
\end{equation}
The constant remains independent of $k$ and $D$.

\subsection{Exact slopes with a common head at each level}
\label{dir:deterministic-heads}
For fixed finite ancestral $A_0\subseteq\cdots\subseteq A_m$, with
$E_i=A_i\setminus A_{i-1}$, the deterministic-head leakage estimate
\cite[Theorem~6.5]{TreePotentials} is
\begin{equation}\label{dir:leakage-input}
 \sum_{i=1}^m\mathbb E\|P_{E_i}W_i\|^2
 \le1024\left(\mathbb E\|W_m\|^2+
                 \sum_{i=1}^m\mathbb E\|P_{E_i}W_{i-1}\|^2\right).
\end{equation}
It imposes no support restriction on the martingale.

Return to the noncontracting normalization and keep the true slopes
$W_i$, so $\|W_i\|\le D$.  Put
$\varepsilon=\min\{1/8,1/\sqrt{k}\}$.  Once the finitely many slopes at
level $i$ have been selected, choose a single finite ancestral $A_i$
containing $A_{i-1}$ when $i\ge1$ and satisfying
\[
 \|(I-P_{A_i})W_i\|\le\varepsilon
\]
for every history at that level.  At each current edge the first-child
slopes are pairwise $2$-separated.  Two have head projections less than
$1/2$ apart; one therefore has tail norm at least $3/4$.  Keep that
middle vertex and its two half-edges as the next fair choice.  After all
choices at this level, enlarge the common head to approximate all next
slopes as well.

For $E_{i+1}=A_{i+1}\setminus A_i$, the first-child outcome satisfies
$\|P_{E_{i+1}}W_{i+1}\|\ge3/4-\varepsilon\ge1/2$.  This outcome has
conditional probability one half.  On all outcomes the parent layer
norm is at most $\varepsilon$, because
$P_{E_{i+1}}W_i=P_{A_{i+1}}(I-P_{A_i})W_i$.
The sets $A_i$ are fixed across each level, exactly as required by
\eqref{dir:leakage-input}.  It follows that
\begin{equation}\label{dir:leakage-conclusion}
 \frac{k}{8}\le1024(D^2+k\varepsilon^2)
 \le1024(D^2+1).
\end{equation}

\subsection{Supported approximations with predictable heads}
\label{dir:predictable-application}
The predictable-layer estimate \cite[Theorem~6.1]{TreePotentials}
applies when $P_0$ is a fixed finite ancestral set,
$P_i$ is finite ancestral and $\mathcal F_{i-1}$-measurable for $i\ge1$,
$P_{i-1}\subseteq P_i$, and $W_i=P_{P_i}W_i$.  If
$\|W_m\|\le C$ on every outcome, it gives
\begin{equation}\label{dir:predictable-input}
 \sum_{i=1}^m\mathbb E\|P_{P_i\setminus P_{i-1}}W_i\|^2\le495C^2.
\end{equation}

For the estimate with constant $495$, keep the noncontracting
normalization.  Set $\alpha=1/[8(k+1)]$.  Approximate the initial true
slope within $\alpha$ by a finite vector $W_0$ supported on a finite
ancestral $P_0$.  At a history of length $i-1$, choose from the true
edge a deviation $z$ with
\[
 \|(I-P_{P_{i-1}})z\|\ge3/4.
\]
To do so, take two of the $2$-separated deviations with projected
difference less than $1/2$, and use the triangle inequality for their
tails.  Approximate $z$ within $\alpha$ by a finite vector $v_i$ and
choose a finite ancestral $P_i\supseteq P_{i-1}$ containing its support.
For the next fair sign put $W_i=W_{i-1}\pm v_i$ and follow the
corresponding actual half-edge.

The sets $P_i$ and vectors $v_i$ are fixed before this sign.  Consequently
the $W_i$ form an exact dyadic martingale, are supported in $P_i$, and
differ from their true slopes by less than $(i+1)\alpha$.  In particular,
$\|W_k\|\le D+1$.  Since the parent has no coordinates outside
$P_{i-1}$, on either sign
\[
 \|P_{P_i\setminus P_{i-1}}W_i\|
 =\|(I-P_{P_{i-1}})v_i\|
 \ge3/4-2\alpha\ge1/2.
\]
The factor two uses only the bound for $I-P_{P_{i-1}}$.
Every hypothesis of \eqref{dir:predictable-input} now holds, giving
\begin{equation}\label{dir:495-conclusion}
 \frac{k}{4}\le495(D+1)^2.
\end{equation}

Under the same predictable support hypotheses, the separate
variation-and-mass-removal estimate
\cite[Theorem~6.3]{TreePotentials} gives
\begin{equation}\label{dir:variation-input}
 \sum_{i=1}^m\mathbb E\|P_{P_i\setminus P_{i-1}}W_i\|^2\le252C^2.
\end{equation}
In particular its bound is $1008$ when $C=2$.

The variation estimate has a useful different graph normalization.
Rescale the graph to upper constant one and lower constant $1/D$, and
put $\varepsilon=\min\{1/(4D),1/(k+1)\}$.  Approximate the initial
slope within $\varepsilon$ as before.  Given a finite head space
$E=\operatorname{ran}P_{P_{i-1}}$, some deviation has distance from $E$
greater than $1/(2D)$.  Otherwise all deviations would have bounded
approximants in $E$ with errors less than $3/(4D)$, and two close
approximants would contradict their separation $2/D$.
Choose this deviation, approximate it by a finite $v_i$ with error less
than $\varepsilon$, and update $W_i=W_{i-1}\pm v_i$ using a predictable
ancestral enlargement $P_i$.  The cumulative error is less than
$(i+1)\varepsilon$, and hence $\|W_k\|\le2$.  Moreover,
\[
 \|P_{P_i\setminus P_{i-1}}W_i\|
 =\|v_i-P_{P_{i-1}}v_i\|
 \ge\operatorname{dist}(v_i,E)>\frac1{4D}.
\]
Using \eqref{dir:variation-input} with $C=2$ proves
\begin{equation}\label{dir:1008-conclusion}
 \frac{k}{16D^2}\le1008.
\end{equation}
The distance-to-head argument avoids requiring the tail projection to
be contractive.

\subsection{A telescoping squared-gap argument}\label{dir:gap-application}
For finite ancestral $A\subseteq A'$, $x=P_Ax$, and $y=P_{A'}y$, the
fresh-layer squared-gap estimate
\cite[Theorem~5.11]{TreePotentials} is
\begin{equation}\label{dir:gap-input}
 \|P_{A'\setminus A}y\|^2\le
 400\left(\frac{\|x+y\|^2+\|x-y\|^2}{2}-\|x\|^2\right).
\end{equation}
The increment $y$ may have a nonzero part in $A$. Averaging a convexity
gain and telescoping is the martingale method for uniformly convex norms
\cite[Theorem~4.6 and its proof]{Pisier2011}. Here the estimate controls
only the part of the increment in the new coordinate layer.

Normalize the graph to upper constant one and lower constant $1/D$.
Let $a=2/D$ and choose
$0<\gamma<\min\{a/8,1/(k+1)\}$.  At the root use a finite ancestral
projection $W_0=P_{A_0}x_0$ of the true slope, with error at most
$\gamma$.  Work successively by levels.  Given a current true slope
and $A_{i-1}$, its $a$-separated deviations have one member $z$ with
$\|(I-P_{A_{i-1}})z\|\ge a/4$, by the same finite-head compactness
argument.  Choose one at each of the finitely many histories.  Enlarge
to a common finite ancestral $A_i$ so that
$\|z-P_{A_i}z\|\le\gamma$ for all these choices.  Set
$v_i=P_{A_i}z$ and define $W_i=W_{i-1}\pm v_i$, associated with the
actual half-edges.  The parent is supported in $A_{i-1}$, the increment
is supported in $A_i$, and every approximation error is at most
$(i+1)\gamma$.  Thus $\|W_i\|\le2$.

For $E_i=A_i\setminus A_{i-1}$, nesting of the projections gives
\[
 P_{E_i}v_i=(P_{A_i}-P_{A_{i-1}})z,
 \qquad \|P_{E_i}v_i\|\ge a/4-\gamma>a/8.
\]
Here the error is exactly $z-P_{A_i}z$, whose norm is at most
$\gamma$; no layer-projection bound enters this comparison.
Apply \eqref{dir:gap-input} at every parent with $x=W_{i-1}$ and
$y=v_i$.  Averaging over a level and then summing the squared-norm
increases yields
\begin{equation}\label{dir:400-conclusion}
 4\ge\mathbb E\|W_k\|^2
 \ge\|W_0\|^2+\frac{k(a/8)^2}{400}
 \ge\frac{k(a/8)^2}{400}.
\end{equation}
Thus this direct argument also excludes a uniform distortion bound.
Only a finite binary tree was averaged after the coordinate choices;
the availability of infinitely many middle vertices was used at each
selection step.
